\section{Exceptional-set transfer at the two-fifteenths scale}
\label{sec:exceptional-transfer}

This section isolates the new analytic interface.  We first state the
short-interval input in the precise form from which we need it.

\begin{theorem}[Guth--Maynard short-interval inputs]
\label{thm:GM-input}
Fix $\rho>0$.  If
$Y\in[X^{2/15+\rho},X^{0.99}]$, then all but (a set that may depend on
$Y$)
\begin{equation}
 O_\rho\!\left(X\exp\{-c_\rho(\log X)^{1/4}\}\right)
 \label{eq:GM-exception-size}
\end{equation}
integers $x\in[X,2X]$ satisfy
\begin{equation}
 \pi(x+Y)-\pi(x)
 =\frac{Y}{\log x}
 +O_\rho\!\left(Y\exp\{-c_\rho(\log X)^{1/4}\}\right).
 \label{eq:GM-prime-count}
\end{equation}
If instead
\[
 Y\in[X^{17/30+\rho},X^{0.99}],
\]
then \eqref{eq:GM-prime-count} holds for every integer
$x\in[X,2X]$.
\end{theorem}

These are Corollaries~1.3 and~1.4 of arXiv:2405.20552v2
\citep{GuthMaynard2026}, after harmless changes in the positive constant
in the stretched exponential.  The almost-all theorem is stated for
integer centers, which is exactly what is needed here.

\subsection{Smooth target moments outside a sparse set}

Let $G\in C_c^1((0,\infty))$ be fixed.  Choose
$0<\gamma<\Gamma$ with $\supp G\subset[\gamma,\Gamma]$.

\begin{lemma}[Smooth translated moments]
\label{lem:smooth-almost-all-moments}
Fix $0<\eps<13/15$.  Uniformly for
\begin{equation}
 X^{2/15+\eps}\leq H\leq X,
 \label{eq:short-H-range}
\end{equation}
 there is a set of integers
 $\mathfrak E=\mathfrak E(X,H,G)\subset[\alpha X,\beta X]$ with
\begin{equation}
 |\mathfrak E|
 \ll_{\eps,W,G}X
 \exp\{-c_{\eps,G}(\log X)^{1/4}\}
 \label{eq:smooth-exception-size}
\end{equation}
such that every $r\in[\alpha X,\beta X]\setminus\mathfrak E$ satisfies,
for every $A>0$,
\begin{align}
 \sum_hG(h/H)\Lambda(r+h)
 &=HI_0(G)+O_{A,\eps,G}\bigl(H(\log X)^{-A}\bigr),
 \label{eq:smooth-Lambda-moment}\\
 \sum_hG(h/H)\Lambda(r+h)^2
 &=HI_0(G)\log X+O_{\eps,G}(H),
 \label{eq:smooth-Lambda-square-moment}\\
 \sum_hG(h/H)(\Lambda(r+h)-1)^2
 &=HI_0(G)\log X+O_{\eps,G}(H).
 \label{eq:smooth-eta-square-moment}
\end{align}
\end{lemma}

\begin{proof}
 It is enough to treat $r$ in one of finitely many dyadic intervals
 contained in $[\alpha X,\beta X]$.  We separate two ranges.

 First suppose $H\leq X^{0.98}$.  Approximate $G$ by a step function on
 \[
  L=\left\lceil\exp\{(\log X)^{1/8}\}\right\rceil
 \]
 equal subintervals of $[\gamma,\Gamma]$.  The prime count
 on each step is a difference of two cumulative counts with lengths
 $y_jH$, where $\gamma\leq y_j\leq\Gamma$.  Apply the almost-all part of
 \cref{thm:GM-input}, with its parameter smaller than $\eps$, to these
 $O(L)$ lengths and take the union of their exceptional sets.  For large
 $X$, the fixed endpoint factors put all these lengths between
 $X^{2/15+\eps/2}$ and $X^{0.99}$.  Since
 $\log L=(\log X)^{1/8}=o((\log X)^{1/4})$, the union does not change the
 shape of \eqref{eq:smooth-exception-size}.  Partial summation gives the
 prime part of \eqref{eq:smooth-Lambda-moment}; the step-approximation
 error is $O(H/L)$ and is therefore smaller than
 $H(\log X)^{-A}$ for every fixed $A$, using one exceptional set
 independent of $A$.

 Now suppose $H>X^{0.98}$.  Put
 $N=\lfloor(\Gamma-\gamma)H/X^{0.97}\rfloor$ and partition
 $[\gamma H,\Gamma H]$ into $N$ equal consecutive blocks.  For large
 $X$, every block length lies between $X^{0.97}$ and $2X^{0.97}$, so
 there is no short remainder; changing integer endpoints by $O(1)$ is
 harmless.
 Since $H\leq X$, every block starts at height comparable to
 $X$.  The all-interval part of \cref{thm:GM-input}, with any fixed small
 value of $\rho$, applies to every block for large $X$.  Freeze $G(h/H)$
 on each block.  The total error from this freezing is
 $O_G(X^{0.97})=O_{A,G}(H(\log X)^{-A})$ for every fixed $A$, while the
 short-interval errors sum to
 $O(H\exp\{-c(\log X)^{1/4}\})$.  Thus the prime part of
 \eqref{eq:smooth-Lambda-moment} holds at every center in this range; we
 may take $\mathfrak E=\varnothing$.

 In both ranges, the omitted proper prime powers contribute at most
 \[
  \ll_G\left(HX^{-1/2}+\log(2X)\right)\log(2X)
  =o_{A,\eps}\bigl(H(\log X)^{-A}\bigr)
 \]
 for every fixed $A$.  This completes the proof of
 \eqref{eq:smooth-Lambda-moment} for the full von Mangoldt function.

For the square, first restrict to prime target values.  On the support,
$r+h\asymp X$, and
\[
 \Lambda(p)^2=(\log p)^2
 =\log X\,\Lambda(p)+O_G(\Lambda(p)).
\]
Equations \eqref{eq:smooth-Lambda-moment} and partial summation therefore
give the main term in \eqref{eq:smooth-Lambda-square-moment}, with an
$O(H)$ remainder.  Proper prime powers in an interval of length $O(H)$
at height $X$ contribute
\[
 \ll_G\left(HX^{-1/2}+\log(2X)\right)(\log(2X))^2=o_\eps(H).
\]
Indeed, for each exponent $j\geq2$, consecutive $j$th powers at height
$X$ are separated by $\gg X^{1-1/j}$, and there are $O(\log X)$ possible
exponents.  This proves \eqref{eq:smooth-Lambda-square-moment}.
Finally expand $(\Lambda-1)^2=\Lambda^2-2\Lambda+1$ and use
$\sum_hG(h/H)=HI_0(G)+O_G(1)$.
\end{proof}

\subsection{The weighted transfer}

The preceding exceptional set could depend on $H$ and on the step
approximation.  Its exact location is immaterial because its mass under
the factor-product weights is negligible.

\begin{proposition}[Weighted exceptional-set transfer]
\label{prop:weighted-exception-transfer}
Let $\mathfrak E\subset[\alpha X,\beta X]\cap\Z$ satisfy
\begin{equation}
 |\mathfrak E|\ll X
 \exp\{-c(\log X)^{1/4}\}
 \label{eq:generic-exception-size}
\end{equation}
for some fixed $c>0$.  Then, uniformly in $D\geq1$ and $t\in\R$, for
every $A>0$,
\begin{align}
 \sum_{r\in\mathfrak E}A_D(r)
 &\ll_{A,c,W}X(\log X)^{-A},
 \label{eq:A-exception-transfer}\\
 \sum_{r\in\mathfrak E}|B_{D,t}(r)|
 &\ll_{A,c,W}X(\log X)^{-A}.
 \label{eq:B-exception-transfer}
\end{align}
\end{proposition}

\begin{proof}
Cauchy--Schwarz and \cref{lem:product-ledgers} give
\[
 \sum_{r\in\mathfrak E}A_D(r)
 \leq|\mathfrak E|^{1/2}
 \left(\sum_rA_D(r)^2\right)^{1/2}
 \ll_WX(\log X)^2e^{-c(\log X)^{1/4}/2}.
\]
The stretched exponential absorbs every fixed power of $\log X$.  The
same argument, using \eqref{eq:B-second-moment}, proves
\eqref{eq:B-exception-transfer}.
\end{proof}

\begin{remark}[What is new and what is imported]
The prime number theorem in almost all short intervals is the deep input
of \citet{GuthMaynard2026}.  The step needed in the present factor-ray
setting is
\cref{prop:weighted-exception-transfer}: it verifies that their
exceptional centers cannot concentrate on the product positions with the
unbounded source weight generated by the factor-ray diagonal.  We do not
claim a new short-interval prime number theorem.
\end{remark}
